\documentclass[a4paper, serif, 10pt]{moderncv}

%% moderncv
% style and color
\moderncvstyle{banking}
\moderncvcolor{black}

%% page layout/margins
\usepackage[top=2.5cm, bottom=2.5cm, left=1.5cm, right=1.5cm]{geometry}

\nopagenumbers{}

%% Babel config for Norwegian language
% ref:
% - https://latex3.github.io/babel/guides/locale-norwegian.html
\usepackage[norsk]{babel}

%% fontspec
\usepackage{fontspec}

\IfFontExistsTF{Linux Libertine}{
  \setmainfont[Ligatures={TeX, Common}, Numbers=OldStyle]{Linux Libertine}
}{}
\IfFontExistsTF{Linux Libertine O}{
  \setmainfont[Ligatures={TeX, Common}, Numbers=OldStyle]{Linux Libertine O}
}{}

%% CTeX
% CJK support, used to show CJK characters in the resume
%
% - fontset=none: disable builtin fontset but instead set the CJK font manually
% - heading=false: disable ctex heading
% - punct=kaiming: use kaiming punctuations styles for CJK
% - scheme=plain: use plain scheme, do not override \normalsize font size
% - space=auto: space settings for CJK characters
%
% ref:
% - http://ctan.mirrorcatalogs.com/language/chinese/ctex/ctex.pdf
\usepackage[UTF8, heading=false, punct=kaiming, scheme=plain, space=auto]{ctex}

\IfFontExistsTF{Noto Serif CJK SC}{
  \setCJKmainfont{Noto Serif CJK SC}
}{}
\IfFontExistsTF{Noto Sans CJK SC}{
  \setCJKsansfont{Noto Sans CJK SC}
}{}

%% line spacing
\usepackage{setspace}
\setstretch{1.125}

%% URL styling - use normal text instead of monospace
\urlstyle{same}

% Auto-underline all links
% Defer until \begin{document} so hyperref is already loaded
\AtBeginDocument{%
  \let\oldhref\href
  \renewcommand{\href}[2]{\oldhref{#1}{\underline{#2}}}%
}

%% Patch \httplink and \httpslink to handle full URLs
% The original moderncv macros always prepend http:// or https:// to URLs.
% This causes issues when the URL already has a protocol (e.g., https://example.com)
% resulting in malformed URLs like https://https://example.com.
% This patch checks if the URL already contains "://" and uses it directly if so.
% Note: We use \str_set:Nx (x-expansion) to expand macros like \@homepage first.
\ExplSyntaxOn
\str_new:N \g_yamlresume_url_str

\RenewDocumentCommand{\httpslink}{O{}m}{
  \str_set:Nx \g_yamlresume_url_str {#2}
  \str_if_in:NnTF \g_yamlresume_url_str {://}
    { \url{#2} }
    { \url{https://#2} }
}

\RenewDocumentCommand{\httplink}{O{}m}{
  \str_set:Nx \g_yamlresume_url_str {#2}
  \str_if_in:NnTF \g_yamlresume_url_str {://}
    { \url{#2} }
    { \url{http://#2} }
}
\ExplSyntaxOff

\name{Ingrid Nordby}{}
\title{DevOps-ingeniør med spesialisering i skyplattformer, automatisering og pålitelighet}
\phone[mobile]{+47 918 45 273}
\email{ingrid.nordby@nordbydrift.no}
\homepage{https://www.ingridnordby.no}

\address{Thorvald Meyers gate 42, Oslo, Oslo, Norge, 0555}{}{}

\extrainfo{{\small \faGithub}\ \href{https://github.com/ingridnordby}{@ingridnordby} {} {} {} • {} {} {} 
{\small \faLinkedin}\ \href{https://www.linkedin.com/in/ingridnordby}{@ingridnordby}}

\begin{document}

\maketitle

\section{Grunnleggende informasjon}

\cvline{}{DevOps-ingeniør med ni års erfaring fra skyleveranser, plattformarbeid og automatisering i \textbf{høyt tilgjengelige} produksjonsmiljøer.
%
\begin{itemize}
\item Bygget og driftet Kubernetes-plattformer som betjener over 200 tjenester og 40 produktteam
\item Automatisert infrastruktur med Terraform, GitOps og policy-som-kode
\item Redusert nedetid og skykostnader gjennom SLO-arbeid, observabilitet og FinOps
\end{itemize}}
\section{Utdanning}

\cventry{aug. 2013–juni 2018}
        {Master, Informatikk, Poeng: 3,8}
        {Norges teknisk-naturvitenskapelige universitet (NTNU)}
        {\url{https://www.ntnu.no}}
        {}
        {\begin{itemize}
\item Fordypning i distribuerte systemer og skalerbar infrastruktur
\item Masteroppgave om automatisk kapasitetsstyring i containerplattformer
\item Gruppeleder i studentforeningen for systemdrift gjennom tre semestre
\end{itemize}
\textbf{Kurs}: Operativsystemer, Distribuerte systemer, Algoritmer og datastrukturer, Databasemodellering, Nettverkssikkerhet, Skyarkitektur og skalerbarhet}

\cventry{aug. 2010–juni 2013}
        {Videregående skole, Informasjonsteknologi, Poeng: 5,4}
        {Bergen tekniske fagskole}
        {\url{https://www.btf.no}}
        {}
        {\textbf{Kurs}: Nettverksteknologi, Linux-administrasjon, Programmering i Python}
\section{Arbeidserfaring}

\cventry{mars 2022–Nå}
        {Lead DevOps-ingeniør}
        {Nordlys Digital AS}
        {\url{https://www.nordlysdigtal.no}}
        {}
        {\begin{itemize}
\item Leder et plattformteam på seks ingeniører med ansvar for Kubernetes-plattformen til alle produktteam
\item Innførte GitOps med Argo CD og Terraform, og kuttet ledetiden for infrastrukturleveranser fra dager til under én time
\item Etablerte SLO-er og feilbudsjetter som løftet tilgjengeligheten fra 99,5 \% til 99,95 \%
\item Reduserte skykostnadene med 31 \% gjennom rettighetsizing, spot-instanser og faste FinOps-rutiner
\item Fungerer som teknisk mentor og driver interne fagdager og onboarding av nye ingeniører
\end{itemize}
\textbf{Nøkkelord}: Kubernetes, Terraform, Argo CD, SRE, FinOps}

\cventry{jan. 2020–mars 2022}
        {Senior DevOps-ingeniør}
        {Skybasert AS}
        {\url{https://www.skybasert.no}}
        {}
        {\begin{itemize}
\item Designet og bygget multi-region AWS-arkitektur med Terraform for en SaaS-plattform med 120 000 aktive brukere
\item Automatiserte CI/CD i GitHub Actions og Jenkins med blågrønne utrullinger og automatisk rollback
\item Bygget observabilitetsstack med Prometheus, Grafana og OpenTelemetry som reduserte MTTR fra 90 til 22 minutter
\item Hardet opp containerbilder med Trivy og OPA-policyer integrert i leveringskjeden
\end{itemize}
\textbf{Nøkkelord}: AWS, GitHub Actions, Prometheus, OpenTelemetry, Sikkerhet}

\cventry{aug. 2018–des. 2019}
        {Driftsingeniør}
        {Fjordkraft Teknologi}
        {\url{https://www.fjordkraftsikt.no}}
        {}
        {\begin{itemize}
\item Driftet Linux-baserte produksjonsmiljøer med Ansible og Puppet for over 60 virtuelle tjenere
\item Migrerte on-prem-tjenester til Azure og automatiserte provisjonering med ARM-maler
\item Etablerte vaktrutiner, runbooks og hendelseshåndtering etter ITIL-prinsipper
\end{itemize}
\textbf{Nøkkelord}: Linux, Ansible, Azure, ITIL}
\section{Språk}

\cvline{Norsk}{Morsmål eller tospråklig nivå}
\cvline{Engelsk}{Fullt profesjonelt nivå \hfill \textbf{Nøkkelord}: Arbeidsspråk i 9 år, TOEFL 108}
\section{Ferdigheter}

\cvline{Skyplattformer}{Ekspert \hfill \textbf{Nøkkelord}: AWS, Azure, Kubernetes, EKS, AKS}
\cvline{Infrastruktur som kode}{Avansert \hfill \textbf{Nøkkelord}: Terraform, Ansible, Pulumi, Helm, CloudFormation}
\cvline{CI/CD og automatisering}{Avansert \hfill \textbf{Nøkkelord}: GitHub Actions, GitLab CI, Argo CD, Jenkins, Tekton}
\cvline{Observabilitet}{Avansert \hfill \textbf{Nøkkelord}: Prometheus, Grafana, OpenTelemetry, Loki, Datadog}
\cvline{Programmering}{Avansert \hfill \textbf{Nøkkelord}: Python, Go, Bash, TypeScript, SQL}
\cvline{Sikkerhet og compliance}{Viderekommen \hfill \textbf{Nøkkelord}: Vault, Trivy, OPA, ISO 27001, Zero Trust}
\section{Utmerkelser}

\cventry{nov. 2023}
        {Nordlys Digital AS}
        {Årets teknologipris}
        {}
        {}
        {Tildelt for arbeidet med selvbetjeningsplattformen som gjorde at produktteamene kunne
provisjonere egne miljøer uten manuell driftsoppfølging.}

\cventry{jan. 2023}
        {Nordic DevOps Community}
        {Beste fagartikkel}
        {}
        {}
        {Prisen ble gitt for artikkelen om plattformteam som produkt og hvordan man måler
verdien av intern utviklererfaring.}
\section{Sertifikater}

\cventry{feb. 2023}
        {AWS}
        {AWS Certified DevOps Engineer – Professional}
        {\url{https://aws.amazon.com/certification/}}
        {}
        {}

\cventry{sep. 2022}
        {HashiCorp}
        {HashiCorp Certified: Terraform Associate}
        {\url{https://www.hashicorp.com/certification/terraform-associate}}
        {}
        {}

\cventry{juni 2021}
        {CNCF}
        {Certified Kubernetes Administrator (CKA)}
        {\url{https://www.cncf.io/certification/cka/}}
        {}
        {}
\section{Publikasjoner}

\cventry{apr. 2023}
        {Plattformteamet som produkt – erfaringer fra norsk næringsliv}
        {Kode24}
        {\url{https://www.kode24.no}}
        {}
        {\begin{itemize}
\item Fagartikkel om hvordan plattformteam kan måle og synliggjøre verdien de skaper for produktteam
\item Beskriver konkrete måltall for ledetid, endringsfrekvens og utviklererfaring
\item Deler erfaringer fra innføring av GitOps og selvbetjening i tre nordiske selskaper
\end{itemize}}
\section{Referanser}

\cventry{\emaillink[marit.solheim@nordlysdigtal.no]{marit.solheim@nordlysdigtal.no}}
        {Faglig leder, Nordlys Digital AS}
        {Marit Solheim}
        {+47 900 12 345}
        {}
        {Ingrid er en av de mest systematiske ingeniørene jeg har jobbet med. Hun kombinerer
dyp teknisk forståelse med evnen til å bygge team og skape varige forbedringer i driften.}
\section{Prosjekter}

\cventry{feb. 2023–nov. 2023}
        {Åpen kildekode-portal som lar utviklere provisjonere miljøer og tjenester selv, med innebygd policykontroll.}
        {Selvbetjeningsportal for utviklere}
        {\url{https://github.com/ingridnordby/selfservice-portal}}
        {}
        {\begin{itemize}
\item Bygget med Backstage og Crossplane på toppen av Kubernetes
\item Reduserte manuelle driftsoppgaver med omtrent 70 \% i pilotteamene
\item Brukes i dag av 25 team i tre nordiske selskaper
\end{itemize}
\textbf{Nøkkelord}: Backstage, Crossplane, Platform Engineering}

\cventry{sep. 2021–aug. 2022}
        {Gjenbrukbare Terraform-moduler for sikre og kostnadseffektive landinger i AWS og Azure.}
        {Terraform-moduler for offentlig sky}
        {\url{https://github.com/ingridnordby/tf-ggv}}
        {}
        {\begin{itemize}
\item Moduler for nettverk, tilgangsstyring, logging og budsjettvarsler
\item Innebygde policyer for kryptering, tagging og nettverkssegmentering
\item Tatt i bruk av ni virksomheter i offentlig sektor
\end{itemize}
\textbf{Nøkkelord}: Terraform, AWS, Azure, Policy as Code}
\section{Interesser}

\cvline{Friluftsliv}{Fjellturer, Langrenn, Padling}
\cvline{Musikk}{Gitar, Kor}
\section{Frivillig arbeid}

\cventry{jan. 2021–jan. 2024}
        {Foredragsholder og fagansvarlig}
        {Nordic DevOps Community}
        {\url{https://www.nordicdevops.no}}
        {}
        {\begin{itemize}
\item Holder foredrag om plattformarbeid, GitOps og observabilitet på møter og konferanser
\item Er med på å organisere den årlige DevOps-dagen for over 400 deltakere
\item Fungerer som mentor for tre junioringeniører gjennom foreningens mentorprogram
\end{itemize}}
    
\end{document}